\documentclass[a4paper,11pt]{article}
\usepackage[T2A]{fontenc}
\usepackage[utf8]{inputenc}
\usepackage[russian]{babel}
\usepackage{amsmath,amssymb,mathtools}
\usepackage{cmbright}
\renewcommand{\familydefault}{\sfdefault}
\usepackage{icomma}
\usepackage[a4paper,left=16mm,right=16mm,top=12mm,bottom=11mm]{geometry}
\usepackage{microtype}
\usepackage{multicol}
\usepackage{xcolor}
\usepackage{hyperref}
\definecolor{accent}{HTML}{15747A}
\definecolor{darkaccent}{HTML}{164C55}
\definecolor{textgray}{HTML}{283238}
\definecolor{muted}{HTML}{65747A}
\definecolor{rulegray}{HTML}{CADADC}
\color{textgray}
\hypersetup{hidelinks}
\setlength{\parindent}{0pt}
\setlength{\parskip}{0pt}
\setlength{\columnsep}{8mm}
\setlength{\abovedisplayskip}{1mm}
\setlength{\belowdisplayskip}{1mm}
\emergencystretch=2em
\newcommand{\sect}[1]{\par\vspace{2.2mm}{\normalsize\bfseries\color{darkaccent}#1}\par\vspace{0.3mm}{\color{rulegray}\rule{\linewidth}{0.45pt}}\par\vspace{0.8mm}}
\newcommand{\fml}[1]{\par\noindent{\normalsize\(\displaystyle #1\)}\par\vspace{2.8mm}}
\begin{document}
\thispagestyle{empty}

\hypersetup{pdftitle={Главные формулы. Ксения Васильченко},pdfauthor={Лёвин Артём Александрович}}
\begin{center}
{\small\bfseries\color{accent}УЧЕБНАЯ ПАМЯТКА \quad\textbullet\quad ОСНОВНЫЕ ФОРМУЛЫ}\par
\vspace{1mm}
{\fontsize{21}{25}\selectfont\bfseries\color{darkaccent} Дроби, делимость и работа\par}
\vspace{0.8mm}
{\small\color{muted}Ксения Васильченко\quad\textbullet\quad 6 класс · математика\quad\textbullet\quad 12.10.2026}
\end{center}
\vspace{2mm}
\noindent{\color{accent}\rule{\linewidth}{0.9pt}}
\vspace{1.3mm}
{\bfseries\color{darkaccent}Главные принципы решения}\par
\vspace{1mm}
\begin{multicols}{2}
\raggedcolumns
\small
\noindent\textbullet\quad Для сложения и сравнения дробей приведите их к общему знаменателю.\par\vspace{0.75mm}
\noindent\textbullet\quad При сокращении делите числитель и знаменатель на общий делитель.\par\vspace{0.75mm}
\columnbreak
\noindent\textbullet\quad Общий знаменатель удобно искать через НОК.\par\vspace{0.75mm}
\noindent\textbullet\quad В задачах на совместную работу складывайте производительности.\par\vspace{0.75mm}
\end{multicols}
\vspace{1.1mm}
\noindent{\color{accent}\rule{\linewidth}{0.5pt}}
\vspace{0.6mm}
\begin{multicols}{2}
\raggedcolumns
\normalsize
\sect{Действия с дробями}
\fml{\frac ab+\frac cd=\frac{ad+bc}{bd}}
\fml{\frac ab-\frac cd=\frac{ad-bc}{bd}}
\fml{\frac ab\cdot\frac cd=\frac{ac}{bd}}
\fml{\frac ab:\frac cd=\frac{ad}{bc}}
\sect{НОД и НОК}
\fml{\operatorname{НОД}(a,b)\cdot\operatorname{НОК}(a,b)=ab}
\fml{\gcd(a,b)=\gcd(b,a\bmod b)}
\columnbreak
\sect{Средние величины}
\fml{\overline x=\frac{x_1+\cdots+x_n}{n}}
\fml{v_{\text{ср}}=\frac{S_{\text{общ}}}{t_{\text{общ}}}}
\sect{Задачи на работу}
\fml{A=pt}
\fml{p_{\text{общ}}=p_1+p_2}
\fml{\frac1T=\frac1{t_1}+\frac1{t_2}\quad(A=1)}
\end{multicols}
\vfill
{\color{accent}\rule{\linewidth}{0.8pt}}\par\vspace{1.5mm}
{\small\color{muted}Лёвин Артём Александрович эксклюзивно для Ксении Васильченко}
\end{document}